% =====================================================================
% Bagian 4 — Results and Discussion
% Paper semifinal Datathon 2026 · format NeurIPS 2015
%
% Setiap angka di bagian ini punya rujukan berkas di docs/. Sel bertanda
% \needdata{} belum punya data dan WAJIB diisi atau dihapus sebelum submit.
% =====================================================================

% \needdata didefinisikan di main.tex

\section{Results and Discussion}

Bagian ini melaporkan seluruh hasil eksperimen, termasuk tiga temuan yang tidak
sesuai hipotesis awal. Pelaporan hasil negatif merupakan komitmen eksplisit pada
kriteria \emph{go/no-go} \emph{concept paper} penyisihan, dan kami
mempertahankannya: menyembunyikan hasil yang merugikan akan merusak validitas
seluruh protokol evaluasi yang justru menjadi kontribusi utama penelitian ini.

\subsection{Partisi data, anti-kebocoran, dan reproduktibilitas}

Pra-pelatihan \emph{encoder} termal memakai 53.011 pasangan RGB--termal FLAME~2
(item~\#9, 254p). Sebelum pelatihan dimulai, 440 indeks bingkai yang menjadi
partisi validasi dan uji RFFNet dikecualikan otomatis dari himpunan latih
FLAME~2; verifikasi anti-kebocoran dijalankan sebagai gerbang keras, dan seluruh
440 bingkai ditemukan padanannya ($\text{missing} = 0$) sehingga tidak ada
bingkai uji yang lolos ke pelatihan. \emph{Fine-tuning} dan evaluasi memakai
anotasi piksel RFFNet dengan partisi resmi repositori, 552 latih, 240 validasi,
200 uji, yang kami verifikasi ulang tidak beririsan sama sekali. Partisi uji
diperlakukan sebagai \emph{held-out} dan dibuka hanya untuk pelaporan akhir.

Sebagai uji reproduktibilitas, kurva degradasi pada partisi validasi dihitung
ulang memakai kerangka evaluasi yang ditulis terpisah dari kerangka aslinya
dengan benih acak yang sama. Keenam titik $\tau$ tereproduksi hingga digit
terakhir, sehingga hasil tidak bergantung pada detail implementasi harness
tertentu.

\subsection{Audit label: satu kelas yang tidak pernah terpelajari}

Kami mengaudit aturan penurunan label dari mask anotasi dan menemukan kekeliruan
pada kerangka evaluasi: predikat keberadaan asap dihitung sebagai dua
pemeriksaan independen atas seluruh citra alih-alih konjungsi per piksel,
sehingga setiap citra yang memuat api otomatis dianggap memuat asap dan kelas
\texttt{fire\_no\_smoke} tidak terjangkau secara konstruksi. Aturan yang
benar, konjungsi per piksel, identik dengan yang dipakai saat
pelatihan, menghasilkan distribusi pada Tabel~\ref{tab:distribusi}.

\begin{table}[!htbp]
\caption{Distribusi kelas setelah koreksi aturan label, dan dampaknya terhadap
akurasi yang dilaporkan.}
\label{tab:distribusi}
\centering
\begin{tabular}{lrrrrr}
\hline
\textbf{Partisi} & \textbf{no\_fire} & \textbf{fire\_no\_smoke} & \textbf{fire\_smoke}
& \textbf{$n$} & \textbf{Pergeseran akurasi} \\
\hline
Latih     & 127 & 1  & 424 & 552 & $-0{,}18$ poin \\
Validasi  & 60  & 4  & 176 & 240 & $-1{,}67$ poin \\
Uji       & 50  & 30 & 120 & 200 & $-15{,}00$ poin \\
\hline
\end{tabular}
\end{table}

Dua konsekuensi harus dinyatakan terbuka. \textbf{Pertama, satu kelas tidak
pernah terpelajari.} Partisi latih hanya memuat satu sampel
\texttt{fire\_no\_smoke}, dan pemeriksaan atas seluruh keluaran model, 17.856
prediksi mencakup tiga konfigurasi modalitas, tiga partisi, dan enam level
degradasi, menunjukkan model \emph{tidak pernah sekali pun} mengeluarkan kelas
tersebut. Karena partisi uji memuat 30 sampel kelas itu, batas atas akurasinya
adalah 85{,}00\%, dan angka tersebut tercapai persis. Sistem ini karenanya
\textbf{tidak dapat membedakan api tanpa asap}, keterbatasan yang justru
menyentuh rezim paling dekat dengan pembakaran membara, sasaran rancangan utama
penelitian ini. \textbf{Kedua, partisi resmi tidak distratifikasi}: proporsi
kelas tersebut adalah 0{,}18\% pada latih, 1{,}67\% pada validasi, dan 15{,}00\%
pada uji. Pergeseran distribusi sebesar ini merupakan properti tolok ukur, bukan
pilihan kami, tetapi membatasi seberapa jauh akurasi pada satu partisi dapat
digeneralisasi ke partisi lain.

Koreksi label ini menggeser akurasi absolut secara seragam pada setiap level
$\tau$ dan setiap konfigurasi modalitas; bentuk kurva degradasi, selisih
antar-konfigurasi, dan seluruh hasil uji signifikansi \textbf{tidak berubah}.

\subsection{Fusi tidak mengungguli unimodal pada kondisi bersih}

Tabel~\ref{tab:bersih} melaporkan hasil pada $\tau = 0$. Hipotesis~\#1
\emph{concept paper} ("fusi $>$ unimodal") \textbf{tidak terkonfirmasi}
terhadap jalur RGB.

\begin{table}[!htbp]
\caption{Akurasi pada kondisi bersih ($\tau = 0$) dengan selang kepercayaan
Wilson 95\%. Uji McNemar berpasangan dua sisi dihitung atas sampel yang sama.}
\label{tab:bersih}
\centering
\begin{tabular}{llrrl}
\hline
\textbf{Partisi} & \textbf{Konfigurasi} & \textbf{Akurasi} & \textbf{CI 95\%}
& \textbf{McNemar vs.\ fusi} \\
\hline
Validasi & Fusi penuh   & 98{,}33\% & [95{,}8; 99{,}4] & (acuan) \\
         & RGB saja     & 98{,}33\% & [95{,}8; 99{,}4] & 0 diskordan; $p = 1{,}000$ \\
         & Termal saja  & 95{,}00\% & [91{,}5; 97{,}1] & 8 diskordan; $p = 0{,}0078$ \\
\hline
Uji      & Fusi penuh   & 85{,}00\% & [79{,}4; 89{,}3] & (acuan) \\
         & RGB saja     & 85{,}00\% & [79{,}4; 89{,}3] & 0 diskordan; $p = 1{,}000$ \\
         & Termal saja  & 85{,}00\% & [79{,}4; 89{,}3] & 0 diskordan; $p = 1{,}000$ \\
\hline
\end{tabular}
\end{table}

Fusi dan RGB-saja bukan sekadar mencapai akurasi yang kebetulan sama: keduanya
menghasilkan \textbf{prediksi yang identik pada setiap sampel} di ketiga
partisi, nol pasangan diskordan. Pada kondisi bersih, membuang kanal termal
tidak mengubah satu pun keputusan model. Sebaliknya, keunggulan fusi atas
termal-saja pada partisi latih dan validasi signifikan secara statistik
($p = 0{,}00098$ dan $p = 0{,}0078$).

Kami menilai temuan ini bukan kegagalan model melainkan properti tolok ukurnya.
Pada partisi uji, ketiga konfigurasi mencapai batas atas yang sama persis, yang
berarti evaluasi kondisi bersih pada tolok ukur ini \emph{tidak memiliki daya
pisah} untuk mendukung klaim arsitektural apa pun. Justru inilah pembenaran
empiris bagi kontribusi metodologis kami: perbedaan antar-konfigurasi hanya
muncul ketika ketahanan diuji secara kontinu.

\subsection{Kurva degradasi: ketahanan, dan satu pembalikan}

Degradasi disintesis dari citra nyata pada kedua modalitas secara serentak
dengan intensitas $\tau \in [0,1]$ yang sama: lapisan asap sintetis
\emph{alpha-blend} termodulasi derau spasial pada kanal RGB, serta
\emph{drift} pola tetap (NUC), derau Gaussian, dan pengaburan pada kanal
termal. Seluruh degradasi diterapkan di ruang piksel mentah sebelum
normalisasi. Kanal termal \textbf{tidak pernah} disintesis dari RGB, batasan
ini dipegang konsisten dengan kritik kami terhadap pendekatan termal sintetis
pada Bagian~2.

\begin{figure}[!htbp]
\centering
\includegraphics[width=0.92\textwidth]{kurva_degradasi_terkoreksi.pdf}
\caption{Kurva degradasi pada partisi validasi dan uji, setelah koreksi label
Subbagian~4.2. Identitas seri disandikan lewat warna, gaya garis, dan penanda
sekaligus sehingga tetap terbaca pada cetakan hitam-putih. Pada partisi
validasi, fusi mempertahankan akurasi tertinggi di seluruh level degradasi.
Pada partisi uji, jalur termal-saja sama sekali tidak terpengaruh, prediksinya
identik di keenam level, sementara fusi tetap lebih tahan daripada RGB-saja.}
\label{fig:kurva}
\end{figure}

Penurunan akurasi dari $\tau = 0$ ke $\tau = 1$, dalam poin persentase untuk
fusi / RGB-saja / termal-saja, adalah $1{,}45$ / $3{,}26$ / $2{,}72$ pada
partisi latih, $\mathbf{1{,}67}$ / $5{,}42$ / $4{,}58$ pada validasi, dan
$\mathbf{4{,}50}$ / $6{,}50$ / $0{,}00$ pada uji.

Satu klaim bertahan konsisten di ketiga partisi: \textbf{fusi lebih tahan
daripada jalur RGB-saja}. Pada validasi fusi mempertahankan akurasi penuh hingga
$\tau = 0{,}6$ sementara RGB-saja sudah turun 2{,}92 poin di titik yang sama.

Namun klaim yang lebih kuat, bahwa fusi adalah konfigurasi paling tahan,
\textbf{tidak bertahan pada partisi uji}. Di sana jalur termal-saja tidak
mengalami penurunan sama sekali. Pemeriksaan pada tingkat sampel menunjukkan
prediksi termal-saja pada partisi uji \emph{identik di keenam level $\tau$};
tidak satu pun keputusan berubah. Pada partisi latih dan validasi, prediksi
jalur yang sama berubah pada 5--17 sampel.

Kami tidak menafsirkan kedataran ini sebagai ketahanan sejati. Penjelasan yang
lebih hemat adalah bahwa margin keputusan jalur termal pada 200 sampel uji
tersebut melampaui apa yang mampu digerus oleh operator degradasi kami,
sehingga yang terukur adalah kelemahan operator, bukan keunggulan modalitas.
Konsekuensi metodologisnya jelas dan kami nyatakan sebagai keterbatasan:
intensitas operator degradasi termal perlu dikalibrasi terhadap margin
keputusan model, bukan ditetapkan \emph{a priori}.

Pembalikan peringkat antar-partisi ini sekaligus memperkuat argumen protokol.
Uji ketahanan biner, membandingkan fusi penuh melawan modalitas tunggal pada
satu partisi, akan melaporkan kesimpulan yang berlawanan tergantung partisi
mana yang kebetulan dipakai. Kurva kontinu pada beberapa partisi memperlihatkan
ketidakstabilan itu, yang justru merupakan informasi.

\subsection{Reliability gating tidak meningkatkan akurasi}

\emph{Gate} keandalan dilatih dengan supervisi regresi terhadap $1 - \tau$,
memakai label degradasi yang dihasilkan sendiri oleh pipeline augmentasi
terkendali, sehingga tidak memerlukan anotasi tambahan. Galat validasi regresi
mencapai $0{,}0860$, dan skor yang dihasilkan berada pada rentang nyata yang
lebar (RGB 15--80\%, termal 22--83\%), bukan terpusat di sekitar nilai
\emph{default}.

Meskipun demikian, \textbf{gating tidak meningkatkan akurasi klasifikasi}. Pada
$\tau \in \{0{,}0;\ 0{,}3;\ 1{,}0\}$ akurasi dengan dan tanpa gating identik,
dan pada $\tau = 0{,}6$ konfigurasi ber-gating lebih rendah satu sampel dari 240
(0{,}42 poin, derau satu sampel). Karena akurasi pada tolok ukur ini nyaris
tidak memiliki ruang gerak, eksperimen ini tidak dapat memisahkan "gating tidak
bermanfaat" dari "tolok ukur tidak mampu mengukurnya", dan kami melaporkan
keduanya sebagai kemungkinan yang belum terpisahkan. Nilai gating yang dapat
kami pertahankan bersifat operasional: skor keandalan per modalitas tampil pada
setiap kartu \emph{alert} sehingga operator mengetahui deteksi tersebut
bersandar pada sensor mana, mengubah \emph{modality dropout} dari regularisasi
pasif menjadi informasi yang sampai ke pengambil keputusan. Angka ablation ini
dihitung sebelum koreksi label Subbagian~4.2; koreksi tersebut menggeser kedua
lengan perbandingan secara seragam sehingga perbandingan antar-lengan tidak
terpengaruh.

\subsection{Lokalisasi dan kesenjangan domain}

Dua jalur lokalisasi dievaluasi pada partisi uji. Jalur A (\emph{attention
rollout}, tanpa label piksel) mencapai \emph{pointing game accuracy} $0{,}6133$
($92/150$ sampel ber-api; CI 95\% $[53{,}3; 68{,}7]$); Jalur B (\emph{head}
segmentasi, dengan label piksel) mencapai mIoU $0{,}5328$. Keduanya
\textbf{tidak diukur pada metrik yang sama}, sehingga ini bukan perbandingan
langsung dan tidak boleh dibaca sebagai salah satu jalur lebih baik. Yang dapat
disimpulkan lebih terbatas namun relevan secara praktis: jalur tanpa label
menunjuk area api pada mayoritas sampel, sehingga ekspansi ke lanskap tanpa
anotasi piksel tetap memungkinkan. Untuk penerapan kami memilih \emph{head}
segmentasi karena mIoU-nya lebih meyakinkan dibanding ambang \emph{pointing
game} yang belum dikalibrasi.

Pada 100 sampel \emph{holdout} FLAME~3 yang tidak beririsan dengan subset
adaptasi, akurasi \emph{zero-shot} mencapai \textbf{92\%} (CI 95\%
$[85{,}0; 95{,}9]$) tanpa adaptasi apa pun. Adaptasi LoRA dengan
$N = 10/25/50/100$ sampel menghasilkan $92\%$, $95\%$, $100\%$, dan $86\%$,
\textbf{tidak monoton}, dan pengulangan menghasilkan pola berbeda karena benih
acak pada inisialisasi bobot dan urutan \emph{batch} tidak dikunci. Kami
karenanya \textbf{tidak} mengklaim hubungan antara jumlah sampel adaptasi dan
performa. Yang dapat dipertahankan hanyalah generalisasi \emph{zero-shot} yang
kuat lintas-dataset, yang menunjukkan kesenjangan domain FLAME~2
$\rightarrow$ FLAME~3 lebih kecil dari dugaan awal untuk tugas ini. Perlu
ditegaskan bahwa keduanya berlokasi di Amerika Serikat dan beriklim sedang,
sehingga angka ini bukan bukti transfer ke gambut tropis.

\subsection{Efisiensi dan kelayakan inferensi tepi}

Model memiliki 30.626.467 parameter (13.896.355 terlatih), menuntut 18{,}68
GFLOPs per \emph{forward}, dan berjalan pada median \textbf{117{,}5 ms per
bingkai} dengan ONNX Runtime fp32 di laptop Lenovo ThinkPad T480. Karena model
beroperasi pada bingkai penuh dan bukan per ubin, angka itu sudah merupakan
latensi per bingkai. Sasaran rancangan $<100$ ms per ubin tercapai pada satu
dari dua konfigurasi CPU yang diuji (Lampiran~\ref{app:efisiensi}).

Dua catatan pembatas. \textbf{Pertama, tidak ada perangkat tepi fisik yang
diuji.} Tidak tersedia akses ke Jetson, Raspberry Pi, maupun ponsel selama
pengerjaan. Angka di atas adalah proksi CPU x86, dan kami tidak mengklaim
validasi pada kelas perangkat mana pun.

\textbf{Kedua, kuantisasi INT8 memperlambat inferensi.} Ukuran model turun
68{,}8\% (117{,}10 $\rightarrow$ 36{,}53 MB), tetapi latensi justru naik menjadi
$0{,}36\times$ kecepatan fp32.
Penjelasan yang paling mungkin adalah biaya konversi tipe data pada lapisan
\emph{attention} Transformer melebihi penghematan aritmetikanya pada CPU. Ini
hasil negatif dan kami laporkan apa adanya. Dampak kuantisasi terhadap akurasi
\textbf{tidak dapat ditetapkan}: evaluasi yang tersedia dijalankan pada partisi
yang dibangun dari rentang bingkai berstatus sementara, dan menghasilkan
metrik degeneratif yang tidak dapat ditafsirkan.

\subsection{Analisis kegagalan}

Seluruh kesalahan pada kondisi bersih di partisi latih dan validasi berasal
dari satu jalur dan satu arah: \textbf{jalur termal-saja menghasilkan alarm
palsu} (11 pada latih, 8 pada validasi), sementara jalur fusi dan RGB tidak
menghasilkan satu pun. Tidak ada kejadian api yang terlewat oleh konfigurasi
mana pun, \emph{recall} deteksi kejadian api bernilai 1{,}000 di seluruh
partisi dan seluruh level $\tau$. Untuk sistem peringatan dini, modalitas yang
tersisa gagal ke arah yang secara operasional lebih murah.

\begin{figure}[!htbp]
\centering
\includegraphics[width=0.56\textwidth]{kasus_gagal_2.png}
\caption{Dua dari empat alarm palsu jalur termal-saja pada partisi validasi.
Label sebenarnya \texttt{no\_fire}, diprediksi \texttt{fire\_smoke}. Keempatnya
memuat permukaan tanah gundul yang terpanaskan matahari, terang luas pada kanal
termal, sementara kanal RGB tidak memperlihatkan nyala maupun asap dan mask
anotasi kosong. Jalur fusi dan RGB tidak melakukan kesalahan pada bingkai-bingkai
ini.}
\label{fig:gagal}
\end{figure}

Gambar~\ref{fig:gagal} memperlihatkan empat kesalahan tersebut. Mekanismenya
konsisten: keempat bingkai memuat \textbf{permukaan tanah gundul yang terpanaskan
matahari}, yang pada kanal termal tampak sebagai area terang luas, sementara
kanal RGB memperlihatkan tidak ada nyala maupun asap. Jalur termal-saja
membaca panas permukaan sebagai api; konteks RGB menyelesaikannya.

Temuan ini memvalidasi secara empiris justifikasi rancangan yang kami tetapkan
\emph{sebelum} eksperimen dijalankan, bahwa konteks RGB berfungsi menekan
\emph{false positive} dari titik panas non-api. Ia juga identik dengan moda
kegagalan yang melekat pada pemicu satelit: FIRMS menandai anomali termal, dan
tungku industri, suar gas, serta tanah panas ikut terdeteksi. Relevansinya
terhadap lahan gambut langsung: permukaan gambut kering yang terbuka memanas
ekstrem pada siang hari, sehingga sistem yang bersandar pada termal saja akan
menghasilkan alarm palsu tepat pada lanskap yang menjadi sasaran rancangan.

Keyakinan model menyertai arah yang benar: pada jalur termal-saja, keyakinan
rata-rata saat keliru adalah $0{,}691$ dibanding $0{,}994$ saat benar. Model
kurang yakin justru ketika salah, yang mendukung pemilihan titik operasi
berbasis ambang keyakinan.

\subsection{Ancaman terhadap validitas}

\textbf{Domain:} seluruh citra berasal dari \emph{prescribed burn} hutan pinus di
Arizona, tanpa satu pun bingkai gambut tropis; sistem dirancang untuk rezim
membara dan divalidasi pada proksi terukur, dan belum tervalidasi untuk gambut.
\textbf{Ukuran sampel dan benih:} seluruh angka berasal dari satu benih acak
tanpa pengulangan, dan pada partisi validasi satu sampel bernilai 0{,}42 poin,
sehingga selisih di bawah satu poin berada dalam batas derau, karenanya setiap
angka utama disertai selang Wilson dan setiap perbandingan memakai McNemar
berpasangan. \textbf{Plafon tolok ukur:} akurasi validasi terbaik tercapai pada
epoch pertama pelatihan, dan ketiga konfigurasi mencapai batas atas yang sama
pada partisi uji. \textbf{Tidak ada baseline eksternal:} perbandingan terhadap
bobot RFFNet terlatih tidak terlaksana dalam jendela waktu kompetisi, sehingga
seluruh perbandingan bersifat internal antar-konfigurasi modalitas dari model
yang sama. \textbf{Operator degradasi tidak terkalibrasi:} intensitasnya
ditetapkan \emph{a priori} dan pada partisi uji terbukti tidak cukup kuat untuk
menggeser satu pun keputusan jalur termal. \textbf{Pra-pelatihan termal lemah:}
\emph{encoder} termal hanya mencapai akurasi validasi $0{,}5148$ pada tahap
pra-pelatihan FLAME~2, yang kami catat sebagai kemungkinan penyebab terbatasnya
kontribusi kanal termal pada kondisi bersih.
